\section{Introduction}
\label{sec:intro}

Knowledge distillation (KD)~\citep{hinton2015distilling} for large language models (LLMs)~\citep{zhao2023survey} has become an essential research direction for enhancing the accessibility and efficiency of Machine-Learning-as-a-Service (MLaaS)~\citep{cai2024llmaas,xu2024survey_kd_llms,yang2024survey_kdllm}. Through knowledge distillation, a small LLM learns to imitate a large teacher LLM's outputs, enabling effective deployment under limited computational or financial resources~\citep{sanh2019distilbert,jiao2020tinybert,wang2020minilm,li2021dynamickd,liang2023less}. Recent studies show that capability distillation, a specialization of knowledge distillation that focuses supervision on a target capability (e.g., instruction-following, reasoning, mathematics, or coding), can substantially improve the student’s performance on downstream tasks that primarily depend on that capability, while achieving these gains at lower serving cost than the large teacher model~\citep{magister2023teaching,shridhar2022distilling,xu2023wizardlm,yue2024distilling}.
% \yd{we seem to mix learning capabilities and knowledge distillation here. They are not interchangable. Let's make their relationship clearer: knowledge distillation does not equal to learning capabilities; learning capabilities is a type of the broader definition of KD.}


Despite rapid progress in capability distillation, most existing methods still treat capabilities as if they can be improved independently~\citep{taori2023stanford, chiang2023vicuna,zhang2024knowledgeable}, even though empirical evidence suggests that optimizing one capability often triggers broad, unintended shifts in the model’s overall capability profile~\cite {zhong2024revisiting, fang2025kddd_survey,cloud2025subliminal}. This mismatch is especially consequential in the budgeted setting: under a fixed token budget, the goal is to maximize downstream task utility, which depends not only on the targeted capability but also on other capabilities the task may implicitly require. If cross-capability interactions are ignored, allocating more tokens to a single target can become inefficient in two ways: the extra improvement on the target can shrink as the budget grows, and the same updates may reduce performance on other tasks that require non-target capabilities. In this case, additional tokens are effectively spent on updates that provide little task-relevant benefit and may even increase deployment risk in MLaaS, which we refer to as \emph{budget waste}. To systematically examine this gap, we take a first step toward measuring and understanding these interactions through a controlled empirical study. Concretely, we repeatedly distill the student under multiple token budgets, each time allocating the entire budget to one target capability drawn from a set of widely recognized core LLM capabilities: General Knowledge (General), Reasoning, Math, Code, Tool use, Long-Context Understanding (LCU), Steerability, and Multilinguality. After each run, we evaluate the resulting student on the full benchmark suite in Table~\ref{tab:capabilities-benchmarks} and record the score change for every capability. This yields a budget-dependent capability transfer matrix in which the diagonal entries capture on-target improvement and the off-diagonal entries capture how training toward one capability redistributes performance across the others. Across budgets, we observe two consistent patterns: \textit{(i) capability-specific distillation induces systematic, budget-dependent transfer to other capabilities rather than isolated improvements}, and \textit{(ii) increasing the budget for a single target produces smaller additional target gains while making the average harm to non-target capabilities more pronounced}, which together explain why naively scaling the budget for one capability can be inefficient.
% \yd{I think you have not clarify any goal yet. what are people trying achieve? you mean the challenges of trying to achieve what? trying to distill one capability without changing the level of other capabilities? or use the least budget to achieve the best distillation for one capability regardless of other capabilities?} 
% First, improving one target capability frequently causes noticeable shifts in other latent capabilities; second, many of the capabilities enhanced during distillation provide little or no benefit to the downstream task, which means that a portion of the optimization budget is spent on capabilities with low task utility. Without understanding these challenges, capability distillation can misallocate training budgets, alter unrelated capabilities, and produce unpredictable behaviors, posing higher cost and safety concerns for MLaaS deployment.

% In this paper, we take a first step toward understanding capability distillation from a more fundamental perspective. Specifically, capability distillation can be interpreted as compressing a large model into a smaller one that retains only the capabilities needed for a target task, making the student a compact representation of the teacher’s task-relevant intelligence. Since classical theories view intelligence as arising from compression \yd{how this claim connects to your story reads vague and suspicious. intelligence cannot arise from any type of compression over anything. Then is this the type of compression your story is referring to?}, where effective behavior is produced by forming compact representations that support prediction and decision-making~\citep{solomonoff1964formal, wallace1999minimum,hutter2000theory}, this perspective suggests that an LLM’s intelligence can be understood as a composition of interacting capabilities that jointly determine its competence. From this viewpoint, two questions naturally follow: \textit{(1) What is the empirical structure of compressed intelligence when expressed through capabilities}, and \textit{(2) How %%\td{consistency with What/what and How/how}
% %how 
% should a limited distillation budget be allocated to acquire the capabilities required for a specific task?} \yd{the two questions are vague. 1 - What can be an empirical structure of compressed intelligence? a star or a triangle? Audience does not know what specifically you are referring to. 2 - budget allocated over what? you already targeted on a capability so we are still allocating budget on other capabilities (if so why? why do not we directly allocate all on the targeted capability?)? or on other data points?} To answer these questions, we conduct extensive experiments to show that: \textit{(i) capability distillation induces systematic, budget-dependent cross-capability transfer}, and \textit{(ii) capability distillation suffers from diminishing returns and negative spillover, leading to substantial budget waste} 
% \yd{still, using private names without definition or explanation wipes out the meaning of this claim.}.\xq{explain in here and in obs 2}

Building on these observed insights, we propose \textit{ReAD}, a \textit{Re}inforcement-guided c\textit{A}pability \textit{D}istillation framework for large language models. Overall, ReAD combines on-the-fly capability-targeted data generation with token-level knowledge distillation, and uses a lightweight contextual bandit to adaptively allocate a fixed budget across interdependent capabilities. Specifically, ReAD first infers a task requirement vector that identifies which capabilities are essential for improving the downstream utility and treats degradations on these dimensions as harmful spillover. It then allocates distillation effort to maximize a proxy reward that favors capability gains aligned with the task requirements while penalizing spillover and budget consumption. In each interval, ReAD samples capability-labeled prompts according to the current allocation, queries the teacher to obtain supervision, updates the student with a standard distillation loss, and uses the resulting capability-profile change to update an uncertainty-aware UCB allocation rule. Extensive experiments demonstrate that ReAD strengthens task-relevant capabilities and reduces wasted budget on low-utility capability updates under budget constraints.

To sum up, this paper makes the following key contributions: 
% \yd{names are not parallel - they do not match each other on the same level. Also, let's make the description more balanced.}

\begin{itemize}[leftmargin=*, itemsep=1pt, topsep=1pt]
    \item \textbf{Empirical study of cross-capability interactions.}
    We show that distilling toward a single capability consistently changes other capabilities, revealing systematic cross-capability interactions that are often overlooked.
    
    \item \textbf{Budgeted capability distillation formulation.}
    We formulate capability distillation as allocating a fixed budget across multiple interacting capabilities, providing objective for improving utility while controlling side effects.
    
    \item \textbf{ReAD: reinforcement-guided capability distillation.}
    ReAD infers task-relevant capabilities, generates capability-targeted distillation data with controllable style and difficulty, and uses an uncertainty-aware contextual bandit to allocate the budget across capabilities.
    
    \item \textbf{Theory analysis of the capability interdependence.}
    We explain cross-capability interactions and diminishing returns in capability distillation.
\end{itemize}